\setcounter{section}{9}

  \zwischenueberschrift{Aufwärmaufgaben}

\inputaufgabe
{}
{

Zeige, dass man jeden
\definitionsverweis {Tangentialvektor}{}{}

\mavergleichskette
{\vergleichskette
{ v
}
{ \in }{ T_PS^2
}
{  }{
}
{    }{
}
{   }{
}
}
{}{}{}
in einem Punkt $P$ auf der
\definitionsverweis {Einheitssphäre}{}{}
durch einen \anfuehrung{uniformen}{}
\definitionsverweis {differenzierbaren Weg}{}{}
auf einem
\definitionsverweis {Großkreis}{}{}
realisieren kann.

}
{} {}

\inputaufgabe
{}
{

Man gebe möglichst einfache Realisierungen für die
\definitionsverweis {Tangentialvektoren}{}{}
in einem Punkt
\mathl{P=(s,t)}{} auf dem Zylindermantel
\mathl{S^1 \times \R}{} an.

}
{} {}

\inputaufgabe
{}
{

Es seien
\mathkor {} {M} {und} {N} {}
\definitionsverweis {differenzierbare Mannigfaltigkeiten}{}{}
und sei
\maabbdisp {\varphi} { M } { N
} {}
eine
\definitionsverweis {differenzierbare Abbildung}{}{.}
Es sei $P \in M$ und
\mathl{Q=\varphi(P)}{} und es seien
\maabbdisp {\gamma_1,\gamma_2} { I } { M
} {}
zwei
\definitionsverweis {differenzierbare Kurven}{}{}
mit einem offenen Intervall
\mathl{0 \in I}{} und
\mathl{\gamma_1 (0) = \gamma_2(0)=P}{.} Es seien
\mathkor {} {\gamma_1} {und} {\gamma_2} {}
im Punkt $P$
\definitionsverweis {tangential äquivalent}{}{.}
Zeige, dass auch die
\definitionsverweis {Verknüpfungen}{}{}
\mathkor {} {\varphi \circ \gamma_1} {und} {\varphi \circ \gamma_2} {}
tangential äquivalent in $Q$ sind.

}
{} {}

\inputaufgabe
{}
{

Man gebe ein Beispiel einer
\definitionsverweis {injektiven}{}{,}
nicht
\definitionsverweis {surjektiven}{}{,}
\definitionsverweis {differenzierbaren Abbildung}{}{}
\maabbdisp {\varphi} {\R} {\R
} {}
derart, dass die zugehörige
\definitionsverweis {Tangentialabbildung}{}{}
\maabbdisp {T_P(\varphi)} {T_P \R} {T_{\varphi(P)} \R
} {}
in jedem Punkt

\mavergleichskette
{\vergleichskette
{ P
}
{ \in }{ \R
}
{  }{
}
{    }{
}
{   }{
}
}
{}{}{}
bijektiv ist.

}
{} {}

\inputaufgabe
{}
{

Man gebe ein Beispiel einer
\definitionsverweis {surjektiven}{}{,}
nicht
\definitionsverweis {injektiven}{}{,}
\definitionsverweis {differenzierbaren Abbildung}{}{}
\maabbdisp {\varphi} {S^1} {S^1
} {}
derart, dass die zugehörige
\definitionsverweis {Tangentialabbildung}{}{}
\maabbdisp {T_P(\varphi)} {T_P S^1} {T_{\varphi(P)} S^1
} {}
in jedem Punkt
\mathl{P \in S^1}{} bijektiv ist.

}
{} {}

\inputaufgabe
{}
{

Zeige, dass auf einer
\definitionsverweis {differenzierbaren Mannigfaltigkeit}{}{} die Addition von Wegen
\maabbdisp {\gamma_1, \gamma_2} { I } { M
} {}
mit
\mathl{\gamma_1(0) =\gamma_2(0)=P \in M}{,} die man durch eine Karte
\maabbdisp {\alpha} { U } { V
} {}
mit
\mathl{\alpha(P)=0}{} aus der Addition im $\R^n$ erhalten kann, im Allgemeinen von der gewählten Karte abhängt.

}
{} {}

\inputaufgabe
{}
{

Zeige, dass die Tangentialabbildung $T_P(\varphi)$ zu
\maabbeledisp {\varphi} {\R^1} {S^1
} { t } {( \cos  t, \sin  t  )
} {,}
in jedem Punkt
\mathl{P \in \R}{} bijektiv ist.

}
{} {}

\inputaufgabe
{}
{

Man gebe ein Beispiel einer
\definitionsverweis {surjektiven}{}{}
\definitionsverweis {differenzierbaren Abbildung}{}{}
\maabbdisp {\varphi} { M } { N
} {}
zwischen zwei
\definitionsverweis {differenzierbaren Mannigfaltigkeiten}{}{}
\mathkor {} {M} {und} {N} {}
derart, dass die zugehörige
\definitionsverweis {Tangentialabbildung}{}{}
\maabbdisp {T_P(\varphi)} {T_PM} {T_{\varphi(P)}N
} {}
in einem Punkt
\mathl{P \in M}{} nicht surjektiv ist.

}
{} {}

\inputaufgabe
{}
{

Man gebe ein Beispiel einer
\definitionsverweis {injektiven}{}{}
\definitionsverweis {differenzierbaren Abbildung}{}{}
\maabbdisp {\varphi} { M } { N
} {}
zwischen zwei
\definitionsverweis {differenzierbaren Mannigfaltigkeiten}{}{}
\mathkor {} {M} {und} {N} {}
derart, dass die zugehörige
\definitionsverweis {Tangentialabbildung}{}{}
\maabbdisp {T_P(\varphi)} {T_PM} {T_{\varphi(P)}N
} {}
in einem Punkt
\mathl{P \in M}{} nicht injektiv ist.

}
{} {}

\inputaufgabe
{}
{

Es sei $M$ eine
\definitionsverweis {differenzierbare Mannigfaltigkeit}{}{}
und

\mavergleichskette
{\vergleichskette
{ P
}
{ \in }{ M
}
{  }{
}
{    }{
}
{   }{
}
}
{}{}{}
ein Punkt. Wir betrachten die folgende Menge.

\mavergleichskettedisp
{\vergleichskette
{ T
}
{ =} { \left\{  (U,f)   \mid   U \subseteq M \text{ offen} , \,  P \in U   , \,  f \in C^1(U,\R)    \right\}
}
{ } {
}
{ } {
}
{ } {
}
}
{}{}{.}
Wir betrachten die
\definitionsverweis {Relation}{}{}

\mathdisp {(U,f) \sim (V,g) : \text{ es gibt eine offene Menge } W \text{ mit } P \in W \subseteq U \cap V \text{ mit } f  \vert_W = g \vert_W} { . }

\aufzaehlungzwei {Zeige, dass dies eine
\definitionsverweis {Äquivalenzrelation}{}{}
auf $T$ ist.
} {Zeige, dass es eine natürliche
\definitionsverweis {Ringstruktur}{}{}
auf der Menge der
\definitionsverweis {Äquivalenzklassen}{}{}
zu dieser Äquivalenzrelation gibt.
}

}
{} {}

\inputaufgabe
{}
{

Es sei $X$ ein
\definitionsverweis {topologischer Raum}{}{,}
$Y$ eine Menge und
\maabbdisp {\varphi} { X } { Y
} {}
eine
\definitionsverweis {Abbildung}{}{.}
Zeige, dass das
\definitionsverweis {Mengensystem}{}{}

\mathdisp {{\mathcal T } = \left\{ V \subseteq Y  \mid  \varphi^{-1}(V) \text{ ist offen in } X        \right\}} {  }

eine
\definitionsverweis {Topologie}{}{}
auf $Y$ definiert, bezüglich der $\varphi$
\definitionsverweis {stetig}{}{}
ist.

}
{} {}

Die in der vorstehenden Aufgabe eingeführte Topologie nennt man \stichwort {Bildtopologie} {.}

\inputaufgabe
{}
{

Zeige, dass auf dem $\R^n$ durch

\mathdisp {P \sim Q, \text{ falls } P-Q \in \Z^n} {  }

eine
\definitionsverweis {Äquivalenzrelation}{}{}
definiert wird. Die
\definitionsverweis {Quotientenmenge}{}{}

\mavergleichskettedisp
{\vergleichskette
{Y
}
{ =} {\R^n /\!\!\sim
}
{ =} {\R^n/\Z^n
}
{ } {
}
{ } {
}
}
{}{}{}
sei mit der
\definitionsverweis {Bildtopologie}{}{}
zur
\definitionsverweis {Quotientenabbildung}{}{}
\maabb {\varphi} {\R^n } { Y
} {}
versehen. Zeige, dass $Y$ ein
\definitionsverweis {Hausdorffraum}{}{}
ist.

}
{} {}

\inputaufgabe
{}
{

Man entwickle die Grundzüge einer Theorie der \anfuehrung{komplexen Mannigfaltigkeiten}{.} Was ist die zugrunde liegende reelle Mannigfaltigkeit, was ist die
\zusatzklammer {komplexe/reelle} {} {}
Dimension, wie sieht der Tangentialraum aus?

}
{} {}

  \zwischenueberschrift{Aufgaben zum Abgeben}

\inputaufgabe
{4}
{

Es sei $M$ eine
\definitionsverweis {differenzierbare Mannigfaltigkeit}{}{} und
\mathl{P \in M}{.} Zeige, dass für eine
\definitionsverweis {differenzierbare Kurve}{}{}
\maabbdisp {\gamma} { I } { M
} {}
mit
\mathl{\gamma(0)= P}{} und
\mathl{a \in \R}{} im
\definitionsverweis {Tangentialraum}{}{}
\mathl{T_PM}{} die Beziehung

\mavergleichskettedisp
{\vergleichskette
{a [\gamma]
}
{ =} { [  \lambda]
}
{ } {
}
{ } {
}
{ } {
}
}
{}{}{} gilt, wobei $\lambda$ durch
\mathl{\lambda(t):= \gamma(at)}{} definiert sei.

}
{} {}

\inputaufgabe
{6}
{

Es sei $M$ eine
$C^k$-\definitionsverweis {Mannigfaltigkeit}{}{}
und

\mavergleichskette
{\vergleichskette
{ P
}
{ \in }{ M
}
{  }{
}
{    }{
}
{   }{
}
}
{}{}{.}
Definiere für
$C^k$-\definitionsverweis {Kurven}{}{}
\maabbdisp {\gamma_1,\gamma_2} { I } { M
} {}
mit

\mavergleichskette
{\vergleichskette
{ \gamma_1(0)
}
{ = }{ \gamma_2(0)
}
{ = }{ P
}
{    }{
}
{   }{
}
}
{}{}{}
eine Äquivalenzrelation, die in einer
\zusatzklammer {jeder} {} {}
\definitionsverweis {Karte}{}{}
die
\definitionsverweis {Ableitungen}{}{}
bis zur Ordnung $k$ berücksichtigt. Wie sehen einfache Vertreter dieser Äquivalenzrelation aus? Definiere eine Vektorraumstruktur auf der
\definitionsverweis {Quotientenmenge}{}{}
und bestimme die
\definitionsverweis {Dimension}{}{.}

}
{} {}

\inputaufgabe
{5}
{

Es sei $M$ eine
\definitionsverweis {differenzierbare Mannigfaltigkeit}{}{}
und
\mathl{P \in M}{.} Wir sagen, dass zwei
\definitionsverweis {Kurven}{}{}
\maabbdisp {\gamma_1, \gamma_2} { I } { M
} {}
mit
\mathl{\gamma_1(0)= \gamma_2(0)= P}{} den gleichen \stichwort {Kurvenkeim} {} definieren, wenn es ein
\mathl{\epsilon >0}{} mit

\mavergleichskettedisp
{\vergleichskette
{ \gamma_1 \!\mid_{ [-\epsilon, \epsilon]}
}
{ =} { \gamma_2 \!\mid_{ [-\epsilon, \epsilon]}
}
{ } {
}
{ } {
}
{ } {
}
}
{}{}{}
gibt.

a) Zeige, dass dies eine
\definitionsverweis {Äquivalenzrelation}{}{}
auf der Menge aller Kurven
\maabb {\gamma} { I } { M
} {}
mit $\gamma(0)= P$
\zusatzklammer {und mit verschiedenen offenen Intervallen
\mathl{0 \in I}{}} {} {}
definiert.

b) Zeige, dass
\definitionsverweis {differenzierbare Kurven}{}{,}
die den gleichen Kurvenkeim repräsentieren, auch den gleichen Tangentialvektor repräsentieren.

}
{} {}

\inputaufgabe
{8}
{

Der
\definitionsverweis {Quotientenraum}{}{}

\mavergleichskette
{\vergleichskette
{ Y
}
{ = }{ \R^n/\Z^n
}
{  }{
}
{    }{
}
{   }{
}
}
{}{}{}
sei mit der
\definitionsverweis {Bildtopologie}{}{}
versehen. Definiere auf $Y$ eine Mannigfaltigkeitsstruktur durch geeignete Karten. Zeige, dass die
\definitionsverweis {Quotientenabbildung}{}{}
\maabbdisp {\varphi} {\R^n } { Y
} {}
eine
\definitionsverweis {differenzierbare Abbildung}{}{}
ist, und dass die
\definitionsverweis {Tangentialabbildung}{}{}
in jedem Punkt ein
\definitionsverweis {Isomorphismus}{}{}
ist.

}
{} {}

 [[Kategorie:Latexseite]]
